\section{วิธีการดำเนินงาน}

% เขียนตามกระบวนการวิทยาการข้อมูล (Data Science Process) 5 ขั้นตอน

โครงงานนี้ดำเนินการตามกระบวนการวิทยาการข้อมูล 5 ขั้นตอน ดังนี้

\subsection{ขั้นการตั้งคำถาม (Question Refinement)}

คำถามที่ต้องการตอบจากข้อมูลที่จะเก็บ มีดังนี้

\begin{enumerate}[label=\arabic*),leftmargin=1.9cm,itemsep=0.15em,topsep=0.2em]
    \item เวลาที่ใช้กับโซเชียลมีเดียต่อวัน สัมพันธ์กับคะแนนสุขภาวะทางจิตอย่างไร และความสัมพันธ์นั้นแรงเพียงใด
    \item การติดตามข่าวเชิงลบ สัมพันธ์กับสุขภาวะทางจิตหรือไม่
    \item นักศึกษาที่เสพเนื้อหาต่างประเภทกัน มีคะแนนสุขภาวะทางจิตแตกต่างกันหรือไม่
    \item เวลาที่ใช้กับโซเชียลมีเดีย สัมพันธ์กับจำนวนชั่วโมงและคุณภาพการนอนอย่างไร
    \item เมื่อควบคุมปัจจัยด้านการนอน การออกกำลังกาย และความกดดันจากการเรียนแล้ว เวลาที่ใช้กับโซเชียลมีเดียยังอธิบายความแตกต่างของสุขภาวะทางจิตได้อยู่หรือไม่
\end{enumerate}

\subsection{ขั้นการเก็บรวบรวมข้อมูล (Data Collection)}

\subsubsection*{เหตุผลในการเลือกข้อมูล (WHY)}

ชุดข้อมูลสาธารณะในหัวข้อนี้ที่ตรวจสอบแล้วเป็นข้อมูลสังเคราะห์
จึงไม่สามารถใช้อธิบายสถานการณ์จริงของนักศึกษาได้
โครงงานนี้จึงเก็บข้อมูลด้วยตนเอง โดยเก็บเฉพาะตัวแปรที่จำเป็นต่อการตอบคำถามในขั้นที่ 1

\subsubsection*{แหล่งข้อมูลและกลุ่มตัวอย่าง (WHO)}

ผู้เก็บและเจ้าของข้อมูลคือคณะผู้จัดทำโครงงาน
ผู้ตอบคือนักศึกษาระดับปริญญาตรี มหาวิทยาลัยขอนแก่น
เลือกกลุ่มตัวอย่างแบบตามสะดวก (convenience sampling)
โดยกระจายลิงก์แบบสอบถามผ่านกลุ่มออนไลน์ของคณะและชั้นปีต่าง ๆ
กลุ่มตัวอย่างเป้าหมายจำนวน 150--200 คน

\subsubsection*{ข้อมูลที่เก็บและชนิดของตัวแปร (WHAT)}

\begin{table}[H]
\centering
\caption{ตัวแปรตาม}
\begin{tabularx}{\textwidth}{|l|L|l|c|}
\hline
\textbf{ตัวแปร} & \textbf{คำอธิบาย} & \textbf{ชนิด} & \textbf{ช่วงค่า} \\
\hline
\texttt{wellbeing\_score} & คะแนนสุขภาวะทางจิตจากแบบวัด WHO-5 & อันตรภาค & 0--100 \\
\hline
\end{tabularx}
\end{table}

\noindent
แบบวัด WHO-5 มี 5 ข้อ แต่ละข้อให้คะแนน 0--5 รวมได้ 0--25 แล้วคูณด้วย 4
เพื่อแปลงเป็นสเกล 0--100 คะแนนยิ่งสูงหมายถึงสุขภาวะทางจิตยิ่งดี

\begin{table}[H]
\centering
\caption{ตัวแปรต้นหลัก}
\begin{tabularx}{\textwidth}{|l|L|l|c|}
\hline
\textbf{ตัวแปร} & \textbf{คำอธิบาย} & \textbf{ชนิด} & \textbf{ช่วงค่า} \\
\hline
\texttt{social\_media\_hours} & ชั่วโมงที่ใช้โซเชียลมีเดียต่อวัน อ่านจาก Screen Time & อัตราส่วน & 0--24 \\
\hline
\texttt{news\_days} & จำนวนวันใน 7 วันที่ผ่านมาที่ติดตามข่าวเชิงลบ & อัตราส่วน & 0--7 \\
\hline
\texttt{content\_type} & ประเภทเนื้อหาที่ดูมากที่สุด & นามบัญญัติ & 8 กลุ่ม \\
\hline
\texttt{main\_platform} & แพลตฟอร์มที่ใช้มากที่สุด & นามบัญญัติ & -- \\
\hline
\end{tabularx}
\end{table}

\begin{table}[H]
\centering
\caption{ตัวแปรควบคุมและตัวแปรลักษณะประชากร}
\begin{tabularx}{\textwidth}{|l|L|l|c|}
\hline
\textbf{ตัวแปร} & \textbf{คำอธิบาย} & \textbf{ชนิด} & \textbf{ช่วงค่า} \\
\hline
\texttt{sleep\_hours} & ชั่วโมงการนอนเฉลี่ยต่อคืน & อัตราส่วน & 0--24 \\
\hline
\texttt{sleep\_quality} & คุณภาพการนอนโดยรวม & เรียงลำดับ & 1--10 \\
\hline
\texttt{exercise\_days} & จำนวนวันที่ออกกำลังกายต่อสัปดาห์ & อัตราส่วน & 0--7 \\
\hline
\texttt{study\_pressure} & ความกดดันจากการเรียน & เรียงลำดับ & 1--10 \\
\hline
\texttt{age} & อายุ & อัตราส่วน & -- \\
\hline
\texttt{gender} & เพศ & นามบัญญัติ & -- \\
\hline
\texttt{year\_of\_study} & ชั้นปี & เรียงลำดับ & 1--4 \\
\hline
\texttt{faculty} & คณะ & นามบัญญัติ & -- \\
\hline
\end{tabularx}
\end{table}

\noindent
ตัวแปรควบคุมทั้งสี่ตัวเก็บเพื่อใช้ควบคุมในสมการถดถอย ไม่ใช่ข้อมูลส่วนเกิน
หากไม่เก็บ อาจสรุปผิดว่าผลที่พบมาจากโซเชียลมีเดีย
ทั้งที่จริงมาจากการนอนน้อยหรือความกดดันจากการเรียน

\subsubsection*{ช่วงเวลาที่เก็บ (WHEN)}

\begin{table}[H]
\centering
\begin{tabular}{|l|l|}
\hline
\textbf{ช่วงเวลา} & \textbf{กิจกรรม} \\
\hline
สัปดาห์ที่ 1 & ออกแบบและร่างแบบสอบถาม \\
\hline
สัปดาห์ที่ 1 & ทดลองใช้กับผู้ตอบ 10 คน แล้วปรับแก้ \\
\hline
สัปดาห์ที่ 2--3 & เปิดเก็บข้อมูลจริงเป็นเวลา 14 วัน \\
\hline
สัปดาห์ที่ 4 & ปิดแบบสอบถามและเริ่มวิเคราะห์ \\
\hline
\end{tabular}
\end{table}

\subsubsection*{วิธีการเก็บข้อมูล (HOW)}

ใช้ Google Forms กระจายลิงก์ผ่านกลุ่มออนไลน์ของนักศึกษา
แบบสอบถามมี 18 ข้อ ใช้เวลาตอบประมาณ 4--5 นาที

จุดสำคัญของวิธีการนี้คือการถามชั่วโมงการใช้งานจากตัวเลขที่บันทึกโดยเครื่อง
ผู้ตอบจะได้รับคำแนะนำให้เปิดดูค่าจริงจากมือถือ
(iPhone: ตั้งค่า $\rightarrow$ Screen Time \quad Android: ตั้งค่า $\rightarrow$ Digital Wellbeing)
แทนการประมาณจากความจำ เพื่อลดความคลาดเคลื่อนจากการรายงานตนเอง

\subsubsection*{โครงสร้างแบบสอบถาม}

\begin{table}[H]
\centering
\begin{tabularx}{\textwidth}{|c|L|l|}
\hline
\textbf{ข้อ} & \textbf{คำถาม} & \textbf{รูปแบบคำตอบ} \\
\hline
-- & คำชี้แจงและการขอความยินยอม & บังคับตอบ \\
\hline
1--4 & อายุ เพศ ชั้นปี คณะ & ตัวเลข/ตัวเลือก \\
\hline
5 & ชั่วโมงใช้โซเชียลมีเดียต่อวัน (จาก Screen Time) & ตัวเลข \\
\hline
6 & แพลตฟอร์มที่ใช้มากที่สุด & ตัวเลือกเดียว \\
\hline
7 & ประเภทเนื้อหาที่ดูมากที่สุด & ตัวเลือกเดียว 8 ข้อ \\
\hline
8 & จำนวนวันใน 7 วันที่ติดตามข่าวหนัก & ตัวเลข 0--7 \\
\hline
9--10 & ชั่วโมงการนอน และคุณภาพการนอน & ตัวเลข / สเกล 1--10 \\
\hline
11--12 & วันที่ออกกำลังกาย และความกดดันจากการเรียน & ตัวเลข / สเกล 1--10 \\
\hline
13--17 & แบบวัดสุขภาวะทางจิต WHO-5 & ตามแบบวัดต้นฉบับ \\
\hline
18 & ข้อดักความใส่ใจ & ตัวเลือกเดียว \\
\hline
\end{tabularx}
\end{table}

\noindent
ตัวเลือกของข้อ 7 ได้แก่ (1) ข่าวสารบ้านเมือง การเมือง เหตุการณ์ปัจจุบัน
(2) บันเทิง เช่น ซีรีส์ อนิเมะ คลิปตลก เพลง
(3) เกมและอีสปอร์ต
(4) กีฬา
(5) ความรู้ การเรียน พัฒนาตนเอง
(6) ชีวิตประจำวันของคนอื่น เช่น เพื่อน คนดัง อินฟลูเอนเซอร์
(7) สินค้า รีวิว ช้อปปิ้ง
(8) อื่น ๆ

เหตุผลที่เลือกแบบวัด WHO-5 แทน PHQ-9 คือ WHO-5 ถามในเชิงบวก
และไม่มีข้อคำถามเกี่ยวกับการทำร้ายตนเอง
จึงเหมาะกับการเก็บข้อมูลโดยนักศึกษาที่ไม่มีบุคลากรทางสุขภาพจิตคอยดูแล
ทั้งนี้จะใช้ข้อความจากฉบับแปลภาษาไทยโดยไม่แก้ไขข้อความ ไม่ตัดข้อ และไม่เปลี่ยนสเกลคำตอบ

\subsection{ขั้นการสำรวจและทำความสะอาดข้อมูล (Data Exploration)}

\subsubsection*{การทำความสะอาดข้อมูล}

\begin{table}[H]
\centering
\begin{tabularx}{\textwidth}{|L|L|}
\hline
\textbf{ปัญหาที่คาดว่าจะพบ} & \textbf{วิธีจัดการ} \\
\hline
ผู้ตอบไม่ผ่านข้อดักความใส่ใจ & ตัดออกทั้งแถว \\
\hline
ตอบเร็วผิดปกติ (น้อยกว่า 90 วินาที) & ตรวจสอบและพิจารณาตัดออก \\
\hline
ค่าที่เป็นไปไม่ได้ เช่น นอน 30 ชั่วโมง & ตรวจสอบและกำหนดเป็นค่าว่าง \\
\hline
ค่าว่างในข้อที่ไม่บังคับตอบ & รายงานสัดส่วน แล้วเลือกระหว่างการเติมด้วยมัธยฐานหรือการตัดแถว \\
\hline
ประเภทเนื้อหาบางกลุ่มมีผู้ตอบน้อยกว่า 15 คน & ยุบเหลือ 4 กลุ่ม ได้แก่ ข่าว บันเทิง ชีวิตคนอื่น และอื่น ๆ \\
\hline
\end{tabularx}
\end{table}

\noindent
หลักเกณฑ์การยุบกลุ่มในแถวสุดท้ายกำหนดไว้ล่วงหน้าก่อนเห็นผลการวิเคราะห์
เพื่อไม่ให้เป็นการเลือกยุบกลุ่มตามผลที่ต้องการ

จากนั้นคำนวณตัวแปรใหม่ \texttt{wellbeing\_score} จากผลรวมคะแนน WHO-5 ทั้ง 5 ข้อ คูณด้วย 4
และตรวจสอบความเชื่อมั่นภายในของแบบวัดด้วยค่าสัมประสิทธิ์แอลฟาของครอนบาค

\subsubsection*{การวิเคราะห์ตัวแปรเดี่ยว (Univariate Analysis)}

รายงานค่าเฉลี่ย มัธยฐาน ส่วนเบี่ยงเบนมาตรฐาน ค่าต่ำสุดและค่าสูงสุดของตัวแปรเชิงปริมาณ
และรายงานความถี่กับร้อยละของตัวแปรเชิงคุณภาพ
แสดงผลด้วยฮิสโทแกรมสำหรับตัวแปรตัวเลข และแผนภูมิแท่งสำหรับตัวแปรเชิงคุณภาพ

\subsubsection*{การวิเคราะห์หลายตัวแปร (Multivariate Analysis)}

\begin{table}[H]
\centering
\begin{tabularx}{\textwidth}{|c|L|L|}
\hline
\textbf{ตอบคำถามข้อ} & \textbf{วิธีวิเคราะห์} & \textbf{กราฟที่ใช้} \\
\hline
1 & สหสัมพันธ์เพียร์สัน รายงานทั้งค่า $r$ และ $r^2$ & Scatter plot พร้อมเส้นแนวโน้ม \\
\hline
2 & สหสัมพันธ์ระหว่างจำนวนวันที่ติดตามข่าวกับคะแนนสุขภาวะ & Scatter plot \\
\hline
3 & เปรียบเทียบค่าเฉลี่ยระหว่างกลุ่มด้วย ANOVA หรือ Kruskal-Wallis & Box plot \\
\hline
4 & สหสัมพันธ์ระหว่างชั่วโมงใช้งานกับการนอน & Scatter plot และ Heatmap \\
\hline
5 & การถดถอยพหุคูณ (ดูขั้นที่ 4) & ตารางสัมประสิทธิ์ \\
\hline
\end{tabularx}
\end{table}

\subsubsection*{การตรวจสอบคุณภาพและที่มาของข้อมูล}

นำข้อมูลที่เก็บเองมาเปรียบเทียบกับชุดข้อมูลสาธารณะที่ตรวจพบว่าเป็นข้อมูลสังเคราะห์
โดยใช้เกณฑ์ต่อไปนี้

\begin{enumerate}[label=\arabic*),leftmargin=1.9cm,itemsep=0.15em,topsep=0.2em]
    \item การกระจายของหลักทศนิยม ข้อมูลจากผู้ตอบจริงจะกระจุกที่เลขกลม ขณะที่ข้อมูลสังเคราะห์กระจายสม่ำเสมอทุกหลัก
    \item สัดส่วนค่าว่าง ข้อมูลจากแบบสอบถามจริงย่อมมีค่าว่าง
    \item ขนาดของค่าสหสัมพันธ์ ชุดข้อมูลสังเคราะห์ที่ตรวจสอบมีค่าสูงถึง 0.80--0.95 ซึ่งไม่พบในข้อมูลแบบสอบถามจริง
\end{enumerate}

\subsection{ขั้นการวิเคราะห์เชิงสถิติ (Modeling)}

ใช้การวิเคราะห์ถดถอยพหุคูณ (Multiple Linear Regression) ตามสมการ

\begin{center}
\texttt{wellbeing\_score $\sim$ social\_media\_hours + news\_days + sleep\_hours}\\
\texttt{+ sleep\_quality + exercise\_days + study\_pressure}
\end{center}

\noindent
เพื่อดูว่าเมื่อควบคุมปัจจัยด้านการนอน การออกกำลังกาย และความกดดันจากการเรียนแล้ว
ชั่วโมงการใช้โซเชียลมีเดียยังอธิบายความแตกต่างของคะแนนสุขภาวะได้อยู่หรือไม่

ค่าที่รายงาน ได้แก่ ค่าสัมประสิทธิ์ของแต่ละตัวแปร ซึ่งตีความได้ว่าเมื่อตัวแปรนั้นเพิ่มขึ้น 1 หน่วย
คะแนนสุขภาวะเปลี่ยนแปลงกี่คะแนนโดยที่ตัวแปรอื่นคงที่ ค่า $p$ ของแต่ละสัมประสิทธิ์
ค่า $R^2$ ของสมการโดยรวม และค่า VIF เพื่อตรวจสอบปัญหาตัวแปรต้นสัมพันธ์กันเอง

นอกจากนี้จะตรวจสอบข้อตกลงเบื้องต้นของการถดถอย ได้แก่ ความเป็นเชิงเส้น
การกระจายของค่าคลาดเคลื่อน และความคงที่ของความแปรปรวน
ด้วยกราฟ residual plot และ Q-Q plot

เครื่องมือที่ใช้คือภาษา Python บน Jupyter Notebook ร่วมกับไลบรารี
pandas, numpy, matplotlib, seaborn, scipy และ statsmodels

\subsection{ขั้นผลลัพธ์ที่คาดหวัง (Expected Outputs)}

\begin{enumerate}[label=\arabic*),leftmargin=1.9cm,itemsep=0.15em,topsep=0.2em]
    \item ค่าสหสัมพันธ์และค่า $r^2$ ระหว่างเวลาที่ใช้กับโซเชียลมีเดียกับคะแนนสุขภาวะทางจิต พร้อมการตีความว่าความสัมพันธ์แรงเพียงใด
    \item ตัวเลขที่อธิบายได้เป็นภาษาที่เข้าใจง่าย เช่น เมื่อเวลาใช้งานเพิ่มขึ้น 1 ชั่วโมงต่อวัน คะแนนสุขภาวะเปลี่ยนแปลงกี่คะแนน
    \item ผลเปรียบเทียบคะแนนสุขภาวะระหว่างกลุ่มที่เสพเนื้อหาต่างประเภทกัน โดยเฉพาะกลุ่มที่ติดตามข่าวเป็นหลัก
    \item ข้อสรุปว่าเมื่อควบคุมปัจจัยอื่นแล้ว ความสัมพันธ์ของโซเชียลมีเดียยังคงอยู่หรือไม่
    \item การเปรียบเทียบผลที่ได้กับงานวิจัยที่มีการเผยแพร่ ว่าสอดคล้องหรือแตกต่างอย่างไร
\end{enumerate}

งานวิจัยขนาดใหญ่ในต่างประเทศจำนวนมากพบว่าความสัมพันธ์ระหว่างการใช้โซเชียลมีเดียกับสุขภาวะทางจิตมีขนาดเล็ก
ด้วยขนาดกลุ่มตัวอย่างประมาณ 200 คน การศึกษานี้จึงมีโอกาสได้ผลว่าไม่พบความสัมพันธ์อย่างมีนัยสำคัญทางสถิติ
กรณีดังกล่าวไม่ถือเป็นความล้มเหลวของโครงงาน แต่เป็นผลการศึกษาที่มีความหมายในตัวเอง
และจะรายงานพร้อมการอภิปรายเปรียบเทียบกับงานวิจัยที่เกี่ยวข้อง
